% Original 15,671-witness anchor. No polynomial or constant is changed.
\subsection{The literal seven-witness semilocal predicate}

For $t\in\mathbb Q$, let $\mathcal T(t)$ mean that there are
$u_0,u_1,u_2,u_3,v_1,v_2,v_3\in\mathbb Q$ with
\[
 u_0^2-2u_1^2-3u_2^2+6u_3^2=1,
 \qquad
 (t-2u_0)^2-8v_1^2-12v_2^2+24v_3^2=4.
\]
Then
\[
 \mathcal T(t)\quad\Longleftrightarrow\quad
 t\in\mathbb Z_{(2)}\cap\mathbb Z_{(3)}.
\]
Indeed, use the quaternion convention $i^2=2$, $j^2=3$ and $ij=-ji$.
Its reduced norm is $x_0^2-2x_1^2-3x_2^2+6x_3^2$ and its reduced trace
is $2x_0$. Put $v_0=(t-2u_0)/2$. The two displayed equations say
exactly that $t$ is a sum of two reduced traces of norm-one elements.
No eighth variable is needed, since $v_0$ is recovered by this rational
expression with fixed nonzero denominator~$2$.

In the notation of \cite[Definition~4 and Proposition~6]{Koenigsmann},
this is $T_{2,3}=S_{2,3}+S_{2,3}$. That proposition identifies
$T_{a,b}$ with the intersection of the local integer rings at the
ramified places, with a separate real interval only when infinity is
ramified. Here the Hilbert symbols are $-1$ at $2$ and $3$, and $1$
elsewhere. At $3$ this follows from $(2/3)=-1$; at every other odd
prime both parameters are units. At $2$, writing $2=2^1\cdot1$ and
$3=2^0\cdot3$, the usual exponent is
$0+1\cdot(3^2-1)/8+0=1$. The algebra splits over $\mathbb R$ because
both parameters are positive. Thus the ramification set is precisely
$\{2,3\}$, which proves the claimed equivalence in both directions.
This uses Koenigsmann's sum of two trace sets, not a predicate with
an additional finite offset.

\subsection{The original anchor and its exact index subgroup}

Write $K=\mathbb Q(\rho)$, $\rho^2=2$, and use the literal normalized
curve and generator
\[
 E:y^2=x^3-(11-6\rho)x,
 \qquad
 P=\left(\frac{57}{8}-\frac72\rho,
          -\frac{165}{8}+\frac{417}{32}\rho\right).
\]
As in the original construction, assume the generator input
$2E(K)=\mathbb ZP$. For a nonzero $S=nP$, put
$f(S)=1/x(S)=f_0+f_1\rho$. Retain the exact constant chain
\[
 c_0=3,\qquad c_{j+1}=c_j^2\quad(0\le j<24),\qquad C=c_{24}.
\]
It forces $C=3^{2^{24}}$, with $25$ shared rational variables and
$25$ quadratic-or-linear equations. The anchor introduces two
variables $a_i$, imposes $Ca_i=f_i$, and imposes $\mathcal T(a_i)$
for each $i\in\{0,1\}$. It has $16$ private witnesses and six
quadratic equations, exactly as in the saved module.

Set $N=2^{24}$ and $m=N/2=2^{23}$. We claim that, for nonzero $n$,
\[
 \text{the anchor holds on }nP
 \quad\Longleftrightarrow\quad
 n\in 2\cdot3^{m-1}\mathbb Z.
\]
First consider the place above~$3$. It is unramified quadratic, with
residue field $\mathbb F_9$ and normalized valuation $v(3)=1$.
For rational $a,b$ one has
$v(a+b\rho)=\min(v_3(a),v_3(b))$: after dividing by the smaller
power of $3$, linear independence of $1,\bar\rho$ over $\mathbb F_3$
precludes cancellation.
The curve has good reduction there. Direct calculation gives
\[
 \bar P=(\bar\rho,0),\qquad
 z(2P)=
 \frac{3396356159088-4135204317756\rho}{3287811625151},
 \qquad z=-x/y.
\]
The reduction of $P$ has order two. The denominator in this expression
is a $3$-unit and both numerator coefficients have valuation one.
Consequently $nP$ is in the formal kernel exactly when $n$ is even,
and $v(z(2P))=1$.

For completeness, the valuation of every such multiple can be computed
without an unproved uniformity assertion. The integral formal group
in the parameter $z$ has inverse $-z$, so its multiplication series
are odd. Their linear coefficient is the corresponding integer.
Thus, for a $3$-unit integer $u$, $[u]z=uz+O(z^3)$ preserves positive
valuation, whereas $[3]z=3z+O(z^3)$ increases valuation by exactly one
at depth at least one. The higher terms have valuation at least
$3v(z)>v(z)+1$. Apply this repeatedly to $[n/2](2P)$, allowing
negative integers by negation. It gives
\[
 v(z(nP))=1+v_3(n),\qquad
 v(f(nP))=2+2v_3(n)\quad(2\mid n).
\]
The second equality follows from
$f=z^2(1-(11-6\rho)f^2)$ on the formal branch.
The integral formal-group construction and its linear multiplication
coefficient are given in \cite[IV\S1, Example~IV.2.2.3 and
Proposition~IV.2.3(a)]{SilvermanAEC2009}. No minimal-model assertion at
the bad prime~$2$ is used below.

At the place above~$2$, normalize $v(2)=1$, so $v(\rho)=1/2$.
The two summands of a rational pair have valuations of opposite
integer parity after multiplication by~$2$, and hence cannot cancel.
The displayed coordinates give
\[
 v(x(P))=-3,\qquad v(y(P))=-9/2,\qquad v(z(P))=3/2.
\]
Thus $P$ lies on the formal branch at infinity, even though reduction
at~$2$ is bad. Every integral multiple lies in this formal subgroup.
One can see the integral branch directly by setting $w=-1/y$:
the equation becomes $w=z^3-(11-6\rho)zw^2$, which has a unique
solution $w\in z^3\mathcal O[[z]]$; its group law is the integral
formal group just cited. In particular $f(nP)$ is integral at~$2$.
Since
\[
 v(a+b\rho)=\min(v_2(a),v_2(b)+1/2),
\]
its rational coefficients are both $2$-integral. As $C$ is a $2$-unit,
the two anchor conditions at~$2$ impose no further restriction on $n$.

At~$3$, the two anchor conditions are equivalent to
$v(f(nP))\ge N$. Positive valuation of $f$ implies negative valuation
of $x$, hence membership in the formal kernel, so $n$ is even.
The preceding exact formula makes the remaining condition
$v_3(n)\ge m-1$. Conversely these two conditions give both required
coordinate divisibilities at~$3$, and the condition at~$2$ already
holds. The proved characterization of $\mathcal T$ supplies its
rational witnesses, proving the subgroup claim.

In particular, anchor indices contain a nonzero finite-index subgroup
and have positive representatives. They can be combined with any
finite collection of requirements that indices be divisible by fixed
integers: take a common multiple with $2\cdot3^{m-1}$. More generally,
any finitely many formal-neighborhood conditions whose preimages are
finite-index subgroups can be imposed in the same way. This statement
does not claim compatibility with arbitrary nonzero residue conditions.

The old and new reciprocal coordinates are equivalent for this anchor.
With $\lambda=3+2\rho$, one has $f_{\rm old}=\lambda^2f$ and
\[
 \begin{pmatrix}f_{{\rm old},0}\\f_{{\rm old},1}\end{pmatrix}
 =\begin{pmatrix}17&24\\12&17\end{pmatrix}
 \begin{pmatrix}f_0\\f_1\end{pmatrix}.
\]
This matrix has determinant one and integral inverse, so it preserves
coordinatewise membership in
$C(\mathbb Z_{(2)}\cap\mathbb Z_{(3)})^2$ in both directions.

\subsection{The unchanged numerical cutoff}

The anchor gives $|n|\ge3^{2^{23}-1}$. Conditional only on the separately
audited original height bound $\log_2H<4783227$, it follows that
$|n|>(20H)^2$. Indeed $\log_2 3>3/2$ and $\log_2 20<5$, and
\[
 3(2^{23}-1)=25165821
 >19132928=4(5+4783227).
\]
Thus the original $24$ squarings suffice whenever that exact height
input is established. This calculation does not replace the required
five-contact height theorem by a six-contact theorem or by a larger
constant. No part of this proof changes the saved polynomial.
